% Part: incompleteness
% Chapter: introduction
% Section: decidability

\documentclass[../../../include/open-logic-section]{subfiles}

\begin{document}

\olfileid{inc}{int}{dec}

\olsection{ناپرېکړتيا او ناتکميلي}

د ناتکميلۍ د قضيو د ګوډلي ثبوت د نحو حسابي کول غواړي۔ خو ان له دې
پرته، يوازې په دې فرض چې يوه تيوري ټولې !!{decidable} اړيکې
!!{represents} کوي، ځينې ښې پايلې تر لاسه کولے شو۔ ثبوت يو قطري
استدلال دے چې د درېدنې د مسئلې د ناپرېکړتيا ثبوت ته ورته دے۔

\begin{thm}
که~$\Gamma$ يوه داسې سازګاره تيوري وي چې هره !!{decidable} اړيکه
!!{represents} کوي، نو~$\Gamma$ !!{decidable} نۀ ده۔
\end{thm}

\begin{proof}
فرض کړئ~$\Gamma$ !!{decidable} وي۔ ښيو چې که~$\Gamma$ هره
!!{decidable} اړيکه !!{represents} کړي، نو بايد ناسازګاره وي۔

!!^{decidable} خاصيتونه (يوځاييزې اړيکې) په هغو !!{formula}s تمثيلېږي
چې يو آزاد متغير لري۔ فرض کړئ~$!A_0(x)$، $!A_1(x)$، \dots، د ټولو
داسې !!{formula}s يوه محاسبه کېدونکې شمېره وي۔ اوس لاندې سټ
$D \subseteq \Nat$ وګورئ:
\[
D = \Setabs{n}{\Gamma \Proves \lnot !A_n(\num{n})}
\]
سټ~$D$ !!{decidable} دے، ځکه دا ازمويلے شو چې~$n \in D$ دے که نۀ:
لومړے~$!A_n(x)$، او له هغې څخه~$\lnot !A_n(\num{n})$ محاسبه کوو۔ په
څرګند ډول په~$!A_n(x)$ کښې د~$x$ د هر آزاد پېښېدنې پر ځاے د
اصطلاح~$\num{n}$ ايښوول
او د~$!A_n(\num{n})$ مخې ته~$\lnot$ ليکل يو ميخانيکي کار دے۔ د فرض
له مخې~$\Gamma$ !!{decidable} ده، نو دا ازمويلے شو چې
$\lnot !A_n(\num{n}) \in \Gamma$ ده که نۀ۔ که وي، $n \in D$؛ او که
نۀ وي، $n \notin D$۔ نو~$D$ هم !!{decidable} دے۔
\textbf{د ژباړې د نسخې سمون (OLCMP-072):} سرچينې د قطري جوړونې په
دوو پرله‌پسې ګامونو کښې د فورمول شاخص غورځولے و؛ دلته په دواړو
ځايونو کښې شاخص بېرته راوړل شوے، څو ازموينه د قطري سټ له تعريف سره
هماغه فورمول وکاروي۔

څرنګه چې~$\Gamma$ ټول !!{decidable} خاصيتونه !!{represents} کوي، نو~$D$ هم
!!{represents} کوي۔ او هغه !!{formula}s چې~$D$ په~$\Gamma$ کښې
!!{represents} کوي ټول د~$!A_0(x)$، $!A_1(x)$، \dots، په لړۍ کښې
دي۔ نو داسې عدد~$d$ ونيسئ چې~$!A_d(x)$ په~$\Gamma$ کښې~$D$
!!{represents} کوي۔ که~$d \notin D$ وي، نو څرنګه چې~$!A_d(x)$
!!{represents}~$D$، $\Gamma \Proves \lnot !A_d(\num{d})$۔ خو دا په دې معنا ده
چې~$d$ د~$D$ ټاکونکی شرط پوره کوي، او له دې امله~$d \in D$۔ دا له
$d \notin D$ سره تناقض لري۔ نو د غيرمستقيم ثبوت له مخې~$d \in D$۔
\textbf{د ژباړې د نسخې سمون (OLCMP-073):} سرچينې په دې واحد کښې د
تمثيل د اصطلاحي نښې په پای کښې زائده انګرېزي~s څلور ځله لرله؛ دلته
ټولې پېښې له زائدې نښې پرته د ټاکل شوې اصطلاح په توګه راوړل شوې دي۔

څرنګه چې~$d \in D$، د~$D$ د تعريف له مخې
$\Gamma \Proves \lnot !A_d(\num{d})$۔ له بلې خوا، څرنګه چې~$!A_d(x)$
په~$\Gamma$ کښې~$D$ !!{represents} کوي،
$\Gamma \Proves !A_d(\num{d})$۔ نو~$\Gamma$ ناسازګاره ده۔
\end{proof}

\begin{explain}
پاسنۍ قضيه ښيي چې هېڅ سازګاره تيوري چې ټولې !!{decidable} اړيکې
!!{represents} کوي، !!{decidable} کېدے نۀ شي۔ موږ به وښيو چې~$\Th{Q}$
ټولې !!{decidable} اړيکې !!{represents} کوي؛ معنا ئې دا ده چې ټولې
هغه تيورۍ چې~$\Th{Q}$ رانغاړي، لکه~$\Th{PA}$ او~$\Th{TA}$، هم دا
کار کوي او له دې امله هغوى هم !!{decidable} نۀ دي۔ (څرنګه چې دا ټولې
تيورۍ په معياري مدل کښې رښتونې دي، ټولې سازګارې دي۔)

دا پايله د ناتکميلۍ د لومړۍ قضيې د يوې کمزورې بڼې د ترلاسه کولو
دپاره هم کارولے شو۔ هره هغه تيوري چې !!{axiomatizable} او
!!{complete} وي، !!{decidable} ده۔ نو سازګارې تيورۍ چې
!!{axiomatizable} وي او ټول !!{decidable} خاصيتونه !!{represents}
کوي، !!{complete} کېدے نۀ شي۔
\end{explain}

\begin{thm}
که~$\Gamma$ !!{axiomatizable} او !!{complete} وي، نو !!{decidable}
ده۔
\end{thm}

\begin{proof}
هره ناسازګاره تيوري !!{decidable} ده، ځکه ناسازګارې تيورۍ ټول
!!{sentence}s رانغاړي؛ نو د «آيا~$!A \in \Gamma$؟» پوښتنې ځواب تل
«هو» وي، يعنې پرېکړه ئې کېدے شي۔

نو فرض کړئ~$\Gamma$ سازګاره، او ورسره !!{axiomatizable} او
!!{complete} وي۔ څرنګه چې~$\Gamma$ !!{axiomatizable} ده، نو
!!{computably enumerable} ده۔ ځکه له~$\Gamma$ د بديهي اصولو څخه ټول
سم !!{derivation}s په يوه محاسبه کېدونکې تابع شمېرلے شو۔ له يوه سم
!!{derivation} څخه هغه !!{sentence} محاسبه کولے شو چې دا ئې
!!{derive}s؛ نو په ګډه داسې محاسبه کېدونکې تابع شته چې د~$\Gamma$ ټولې
قضيې شمېري۔ څرنګه چې~$\Gamma$ سازګاره او !!{complete} ده،
!!{sentence} هله او يوازې هله د~$\Gamma$ قضيه ده چې~$\lnot !A$
قضيه نۀ وي۔ نو داسې پرېکړه کولے شو چې~$!A \in \Gamma$ ده که نۀ۔ د
$\Gamma$ ټولې قضيې وشمېرئ۔ کله چې~$!A$ په لست کښې راشي، پوهېږو چې
$\Gamma \Proves !A$۔ کله چې~$\lnot !A$ په لست کښې راشي، پوهېږو چې
$\Gamma \Proves/ !A$۔ څرنګه چې~$\Gamma$ !!{complete} ده، له دې دوو
حالتونو يو بالاخره رامنځته کېږي؛ نو کړنلاره بالاخره ځواب ورکوي۔
\end{proof}

\begin{cor}
\ollabel{cor:incompleteness}
که~$\Gamma$ سازګاره او !!{axiomatizable} وي او هر !!{decidable}
خاصيت !!{represents} کړي، نو !!{complete} نۀ ده۔
\end{cor}

\begin{proof}
که~$\Gamma$ !!{complete} وے، د تېرې قضيې له مخې به !!{decidable} وه
(ځکه !!{axiomatizable} او سازګاره ده)۔ خو څرنګه چې~$\Gamma$ هر !!{decidable}
خاصيت !!{represents} کوي، د لومړۍ قضيې له مخې !!{decidable} نۀ ده۔
\end{proof}

\begin{prob}
وښيئ چې~$\Th{TA} = \Setabs{!A}{\Sat{N}{!A}}$، !!{axiomatizable} نۀ
ده۔ دا فرضولے شئ چې~$\Th{TA}$ ټول د پرېکړې وړ خاصيتونه تمثيلوي۔
\end{prob}

کله چې ثابته کړو، مثلاً، $\Th{Q}$ ټول !!{decidable} خاصيتونه
!!{represents} کوي، ملازمه راته وايي چې~$\Th{Q}$ بايد ناتکميله وي۔
خو ثبوت ئې د خپلواکې !!{sentence} بېلګه نۀ راکوي؛ يوازې ښيي چې داسې
!!a{sentence} بايد موجوده وي۔ د بېلګې د ترلاسه کولو دپاره بايد نحو
حسابي کړو او د ګوډل د اصلي ثبوت نظر تعقيب کړو۔ او البته لا هم بايد
لومړۍ ادعا ثابته کړو، يعنې دا چې~$\Th{Q}$ په رښتيا ټول !!{decidable}
خاصيتونه !!{represents} کوي۔

بايد پام وشي چې هره \emph{په زړه پورې} تيوري ناتکميله يا ناپرېکړه
کېدونکې نۀ ده۔ ډېرې داسې تيورۍ شته چې د په زړه پورې رياضياتي حقايقو
بيانولو ته کافي پياوړې دي، خو د ګوډل د پايلې شرطونه نۀ پوره کوي۔ د
بېلګې په توګه،
$\Th{Pres} = \Setabs{!A \in \Lang{L_{A^+}}}{\Sat{N}{!A}}$، يعنې د
ضرب~$\times$ پرته د حساب د ژبې د هغو !!{sentence}s سټ چې په معياري
مدل کښې رښتونې دي، هم بشپړه او هم د پرېکړې وړ ده۔ دې تيورۍ ته د
پرېسبورګر حساب ويل کېږي، او د طبيعي عددونو په باب ټول هغه حقايق
ثابتوي چې يوازې په~$\Obj{0}$، $\prime$ او~$+$ جوړېدے شي۔

\end{document}
